\chapter{Product Overview}
\section{SaaS Model}
The platform is designed as a multi-tenant SaaS product. The current tenant context identifies \texttt{VidyaSetu Education Group}, \texttt{Saraswati Public School}, \texttt{Pune Wakad Campus}, and academic year \texttt{2026-27}. These values are mock context values and should later be resolved from authenticated request/session state.

\begin{figure}[htbp]
\centering
\begin{tikzpicture}[node distance=14mm, box/.style={draw,rounded corners,align=center,minimum width=30mm,minimum height=10mm,fill=SurfaceMuted}, arr/.style={-{Latex},thick}]
\node[box] (users) {School Users\\Admin, Principal, Teacher, HR};
\node[box,right=of users,minimum width=42mm] (erp) {SchoolERP OS\\Next.js Application};
\node[box,right=of erp] (services) {Tenant-scoped\\Services};
\node[box,below=of services] (mock) {Mock/In-memory\\Data};
\draw[arr] (users) -- (erp);
\draw[arr] (erp) -- (services);
\draw[arr] (services) -- (mock);
\end{tikzpicture}
\caption{Current system context for the implemented frontend and mock service architecture.}
\label{fig:system-context}
\end{figure}


\section{Major Modules}
\begin{longtable}{@{}p{0.32\textwidth}p{0.18\textwidth}p{0.40\textwidth}@{}}
\toprule
\textbf{Module} & \textbf{Status} & \textbf{Notes}\\
\midrule
Student Management & \statusimplemented & Directory, profile, create/edit, lifecycle, placements, documents.\\
Guardian Management & \statusimplemented & Guardian records and student-family relationship surfaces.\\
Academic Structure & \statusimplemented & Academic years, classes, sections, placements, promotions.\\
Attendance & \statusimplemented & Student attendance workstation, history, reports, integrity, exports, and staff attendance.\\
Communication & \statusimplemented & Inbox, conversations, groups, channels, broadcasts, notices, rich text, attachments.\\
Staff Management & \statusimplemented & Staff phases 0--8 implemented in the current source tree.\\
Fees and Exams & \statuspartial & Navigation and dashboard surfaces exist; full domain implementation is not currently specified.\\
Payroll & \statusnot & Not implemented.\\
\bottomrule
\end{longtable}

\section{Workflow Philosophy}
The implemented UI favors operational workstations rather than marketing pages. Examples include the exceptions-first attendance marking screen, Staff 360, staff workload dashboard, and promotion audit views.
